%------------------------------------------------------------------------------%
\documentclass[fleqn,utf8,aspectratio=169,12pt]{beamer}

\usepackage{integracao_e_series}

\title{Integrais Indefinidas}

%------------------------------------------------------------------------------%
\begin{document}

\addtitle

%------------------------------------------------------------------------------%
\section{Integrais Indefinidas}

%------------------------------------------------------------------------------%
\begin{frame}{Integrais Indefinidas}
  Sinônimos
  \begin{itemize}[<+->]
    \item Integral indefinida
    \item Primitiva
    \item Antiderivada
  \end{itemize}
\end{frame}

%------------------------------------------------------------------------------%
\begin{frame}{Definição}
  Uma função $F$ é a \emph{Primitiva} da função $f$
  \pause
  em um intervalo $I$
  \pause
  se
  \[
    \frac{dF}{dx}(x) = f(x)
  \]
  \pause
  para todo $x\in I$
\end{frame}

%------------------------------------------------------------------------------%
\begin{frame}{Primitiva Mais Geral}
  Se $G$ é uma primitiva de $f$ no intervalo $I$,
  \pause
  então a \emph{primitiva mais geral} de $f$ em $I$
  \pause
  é
  \[
    F(x) = G(x) + C
  \]
  \pause
  onde $C$ é uma \emph{constante arbitrária}
  \pause
  (desconhecida)
\end{frame}

%------------------------------------------------------------------------------%
\begin{frame}{Notação}
  Notação da Primitiva, Antiderivada e Integral Indefinida de uma função $f$

  \[
    \int f(x) \, dx
  \]
\end{frame}

%------------------------------------------------------------------------------%
\begin{frame}{Relação Entre Derivada e Primitiva}
  \[
    F(x) = \int f(x)\,dx
    \qquad \Leftrightarrow \qquad
    \frac{dF}{dx} = f(x)
  \]

  \pause
  \[
    \frac{d}{dx} \int f(x)\,dx \;=\; f(x)
  \]

  \pause
  \[
    \int \frac{df}{dx} \,dx \;=\; f(x)
  \]
\end{frame}

%------------------------------------------------------------------------------%
\section{Encontrando uma Primitiva}

%------------------------------------------------------------------------------%
\begin{frame}{Calculando Inversas -- Raízes Quadradas}
\(\sqrt{x}\) é a operação inversa de \(x^2\) para \(x\geq0\)

\pause
\bigskip
\begin{columns}
\column{0.5\linewidth}
  Sabemos calcular o quadrado
  \begin{align*}
    \vphantom{\sqrt{9}} 1^2 & = 1 \\
    \vphantom{\sqrt{9}} 2^2 & = 4 \\
    \vphantom{\sqrt{9}} 3^2 & = 9 \\
    \vphantom{\sqrt{9}} 4^2 & = 16
  \end{align*}
\column{0.5\linewidth}
  \pause
  Valores conhecidos para raízes
  \begin{align*}
    \sqrt{1}  & = 1 \\
    \sqrt{4}  & = 2 \\
    \sqrt{9}  & = 3 \\
    \sqrt{16} & = 4
  \end{align*}
\end{columns}

\pause
\bigskip
Para os demais valores precisamos de uma nova técnica
\end{frame}

%------------------------------------------------------------------------------%
%------------------------------------------------------------------------------%
\begin{question}
  Encontrar a primitiva mais geral de \(f(x)= \cos(x)\)
\end{question}

%------------------------------------------------------------------------------%
\begin{resolution}{Solução}
  Sabemos que
  \[
    G(x) = \sin(x)
  \]
  é uma \emph{primitiva} para $f$,
  \pause
  pois
  \[
    G'(x)
    \pause
    \;=\; \frac{d}{dx}\sin(x)
    \pause
    \;=\; \cos(x)
    \pause
    \;=\; f(x)
  \]
  \pause
  Portanto, a \emph{primitiva mais geral} de $f$ é
  \pause
  \[
    F(x)
    \pause
    \;=\; \int \cos(x)\,dx
    \pause
    \;=\; G(x) + C
    \pause
    \;=\; \sin(x) + C
  \]
\end{resolution}

%------------------------------------------------------------------------------%

%------------------------------------------------------------------------------%
\begin{frame}{Tabela de Primitivas}
\begin{columns}
\column{0.5\linewidth}
  \begin{gather*}
    \onslide<+->{\int x^n         \,dx = \frac{x^{n+1}}{n+1} + c, \quad n \neq -1 \\[1mm]}
    \onslide<+->{\int \frac{1}{x} \,dx = \ln\abs{x} + c \\[1mm]}
    \onslide<+->{\int e^x         \,dx = e^x + c        \\[1mm]}
    \onslide<+->{\int \sin(x)     \,dx = -\cos(x) + c   \\[1mm]}
    \onslide<+->{\int \cos(x)     \,dx = \sin(x) + c}
  \end{gather*}
\column{0.5\linewidth}
  \begin{gather*}
    \onslide<+->{\int \sec^2(x)               \,dx = \tan(x) + c     \\[1mm]}
    \onslide<+->{\int \sec(x) \tan(x)         \,dx = \sec(x) + c     \\[1mm]}
    \onslide<+->{\int \frac{1}{\sqrt{1-x^2}}  \,dx = \arcsin(x) + c  \\[1mm]}
    \onslide<+->{\int \frac{-1}{\sqrt{1-x^2}} \,dx = \arccos(x) + c  \\[1mm]}
    \onslide<+->{\int \frac{1}{1+x^2}         \,dx = \arctan(x) + c}
  \end{gather*}
\end{columns}
\end{frame}

%------------------------------------------------------------------------------%
\section{Propriedades}

%------------------------------------------------------------------------------%
\begin{frame}{Linearidade da Integral Indefinida}
  Sejam $f$ e $g$ integráveis e $k$ uma constante
  \pause
  \bigskip
  \[
    \int f(x) \,\emph{+}\, g(x) \,dx
    \pause
    \;=\; \int f(x) \,dx \,\emph{+} \int g(x) \,dx
  \]
  \pause
  \[
    \int \emph{k} f(x) \,dx
    \pause
    \;=\; \emph{k} \int f(x) \,dx
  \]
\end{frame}

%------------------------------------------------------------------------------%
%------------------------------------------------------------------------------%
\begin{question}
  Encontre a primitiva mais geral da função \(f(x)=3x^2 + 2x + 5\)
\end{question}

%------------------------------------------------------------------------------%
\begin{resolution}{Solução}
  \begin{align*}
    \onslide<+->{F(x)
    & = \int f(x) dx \\}
    \onslide<+->{& = \int 3x^2 + 2x + 5 \,dx \\}
    \onslide<+->{& = \int 3x^2 dx
                   + \int 2x\, dx
                   + \int 5 dx
      & & \text{(Linearidade)}\\}
    \onslide<+->{& = 3\int x^2 dx
                   + 2\int x\, dx
                   + 5 \int dx
      & & \text{(Linearidade)}\\}
    \onslide<+->{& = 3 \frac{x^{2+1}}{2+1}
                   + 2 \frac{x^{1+1}}{1+1}
                   + 5 \frac{x^{0+1}}{0+1} + C
      & & \text{(Tabela)}\\}
    \onslide<+->{& = x^3 + x^2 + 5x + C}
  \end{align*}
\end{resolution}

%------------------------------------------------------------------------------%


%------------------------------------------------------------------------------%
\begin{frame}{Multiplicação da Variável por Constante}
  Qual a Primitiva da função \(f(ax+b)\), onde $a$ e $b$ são constantes?

  \pause
  Se \(F(x)\) é uma primitiva de \(f(x)\)
  \[
    \pause
    \frac{d}{dx}\big(F({\color{magenta}a}x+b)\big)
    \pause
    \;=\; F'({\color{magenta}a}x+b) \, \frac{d}{dx}({\color{magenta}a}x+b)
    \pause
    \;=\; f({\color{magenta}a}x+b) {\color{magenta}a}
  \]
  \[
    \pause
    \frac{d}{dx}\left(\frac{F({\color{magenta}a}x+b)}{{\color{magenta}a}}\right)
    \;=\; f({\color{magenta}a}x+b)
  \]

  \pause
  Então
  \emph{\[
    \int f(ax+b) dx \;=\; \frac{F(ax+b)}{a} + C
  \]}
\end{frame}

%------------------------------------------------------------------------------%
\section{Exemplos}

%------------------------------------------------------------------------------%
\begin{question}
  Calcule \(\int \cos(3\pi x) - \sqrt{x} \, dx\)
\end{question}

\begin{resolution}{Solução}
  \begin{align*}
    \onslide<+->{F(x)}
    \onslide<+->{& = \int \cos(3\pi x) - \sqrt{x} \, dx \\}
    \onslide<+->{& = \int \cos(3\pi x) dx
                   - \int \sqrt{x} \, dx \\}
    \onslide<+->{& = \int \cos(3\pi x) dx
                   - \int x^{\sfrac{1}{2}} \, dx \\}
    \onslide<+->{& = \frac{\sin(3\pi x)}{3\pi}
                   - \frac{x^{\sfrac{1}{2}+1}}{\sfrac{1}{2}+1} + C \\}
    \onslide<+->{& = \frac{\sin(3\pi x)}{3\pi}
                   - \frac{2x^{\sfrac{3}{2}}}{3}
                   + C}
  \end{align*}
\end{resolution}

%------------------------------------------------------------------------------%
%------------------------------------------------------------------------------%
\begin{question}
  Encontre todas as funções $g$, tais que
  \[
    g'(x) = 4 \sin(x) + \frac{2x^5 - \sqrt{x}}{x}
  \]
\end{question}

%------------------------------------------------------------------------------%
\begin{resolution}{Solução}
  \onslide<+->{Queremos as primitivas de $g'$}

  \onslide<+->{Reescrevemos $g'$ como}
  \begin{align*}
    \onslide<+->{g'(x)}
    \onslide<+->{& = 4 \sin(x) + \frac{2x^5- \sqrt{x}}{x} \\}
    \onslide<+->{& = 4 \sin(x) + \frac{2x^5}{x}
                   - \frac{x^{\sfrac{1}{2}}}{x} \\}
    \onslide<+->{& = 4 \sin(x) + 2x^4 -x^{-\sfrac{1}{2}}}
  \end{align*}
\end{resolution}

%------------------------------------------------------------------------------%
\begin{resolution}{Solução}
  \begin{align*}
    \onslide<+->{g(x)}
    \onslide<+->{& = \int g'(x) dx \\}
    \onslide<+->{& = \int 4 \sin(x) + 2x^4 -x^{-\sfrac{1}{2}} dx \\}
    \onslide<+->{& = 4 \int \sin(x) dx
                   + 2 \int x^4 dx
                   - \int x^{-\sfrac{1}{2}} dx \\}
    \onslide<+->{& = 4\big(- \cos(x)\big)
                   + 2 \frac{x^5}{5}
                   - \frac{x^{\sfrac{1}{2}}}{\sfrac{1}{2}}
                   + C \\}
    \onslide<+->{& = -4 \cos(x)
                   + \frac{2}{5}x^5
                   - 2 \sqrt{x}
                   + C}
  \end{align*}
\end{resolution}

%------------------------------------------------------------------------------%
%------------------------------------------------------------------------------%
\begin{question}
  Encontre $f$ tal que
  \[
    f'(x) = e^x + 20\left(1+x^2\right)^{-1}
  \]
  e \(f(0)=-2\)
\end{question}

\begin{resolution}{Solução}
  Primeiro vamos encontrar a família de funções $f$ que são primitivas de $f'$

  \bigskip
  \pause
  Depois usamos a condição \(f(0)=-2\) para encontrar a primitiva solicitada
\end{resolution}

\begin{resolution}{Solução}
  \onslide<+->{Primitiva mais geral}
  \begin{align*}
    \onslide<+->{f(x) & = \int f'(x) dx \\}
    \onslide<+->{& = \int e^x+20\left(1+x^2\right)^{-1} dx \\}
    \onslide<+->{& = \int e^x + 20\frac{1}{1+x^2} dx  \\}
    \onslide<+->{& = \int e^x dx + 20\int \frac{1}{1+x^2} dx  \\}
    \onslide<+->{& = e^x + 20\arctan(x) + C}
  \end{align*}
\end{resolution}

\begin{resolution}{Solução}
  \onslide<+->{Vamos determinar o valor de $C$ usando $f(0)=-2$}
  \begin{align*}
    \onslide<+->{f(0)                    & = -2                       \\}
    \onslide<+->{e^0 + 20 \arctan(0) + C & = -2                       \\}
    \onslide<+->{C                       & = -2 - e^0 - 20 \arctan(0) \\}
    \onslide<+->{                        & = -2 - 1 - 0               \\}
    \onslide<+->{                        & = -3}
  \end{align*}
  \onslide<+->{Portanto
  \[
    f(x) = e^x + 20 \arctan(x) -3
  \]}
\end{resolution}

%------------------------------------------------------------------------------%
%------------------------------------------------------------------------------%
\begin{question}
  Calcule a integral
  \(
    \int x^2\sin(\pi t) \, dt
  \)
\end{question}

%------------------------------------------------------------------------------%
\begin{resolution}{Solução}
  \begin{align*}
    \onslide<+->{F(t) & = \int x^2\sin(\pi t) \, dt}  \\[2mm]
    \onslide<+->{     & = x^2 \int \sin(\pi t) \, dt}  \\[2mm]
    \onslide<+->{     & = x^2 \big(-\cos(\pi t) \big) \frac{1}{\pi} + c} \\[2mm]
    \onslide<+->{     & = -\frac{x^2}{\pi} \cos(\pi t) + c}
  \end{align*}
\end{resolution}

%------------------------------------------------------------------------------%
%------------------------------------------------------------------------------%
\begin{question}
  Calcule a integral
  \(
    \int a(3x-1)^4 - \frac{1}{cx} + \sqrt{x} + 7 \, dx
  \)
\end{question}

%------------------------------------------------------------------------------%
\begin{resolution}{Solução}
\begin{align*}
  \onslide<+->{F(x)
  & = \int a(3x-1)^4 - \frac{1}{cx} + \sqrt{x} + 7 \, dx}
  \\[2mm]
  \onslide<+->{& = a\int (3x-1)^4 \, dx}
  \onslide<+->{- \frac{1}{c}\int \frac{1}{x} \, dx}
  \onslide<+->{  + \int \sqrt{x} \, dx}
  \onslide<+->{  + 7 \int dx}
  \\[2mm]
  \onslide<+->{& = a\int (3x-1)^4 \, dx}
  \onslide<+->{  - \frac{1}{c}\int x^{-1} \, dx}
  \onslide<+->{  + \int x^{\sfrac{1}{2}} \, dx}
  \onslide<+->{  + 7\int x^0 \, dx}
\end{align*}
\end{resolution}

%------------------------------------------------------------------------------%
\begin{resolution}{Solução}
  \begin{align*}
    \onslide<+->{F(x)
    & = a\int (3x-1)^4 \, dx
      - \frac{1}{c}\int x^{-1} \, dx
      + \int x^{\sfrac{1}{2}} \, dx
      + 7\int x^0 \, dx}
    \\
    \onslide<+->{& = a\frac{(3x-1)^{4+1}}{3(4+1)}}
    \onslide<+->{  - \frac{1}{c}\ln\abs{x}}
    \onslide<+->{  + \frac{x^{\sfrac{1}{2}+1}}{\sfrac{1}{2}+1}}
    \onslide<+->{  + 7\frac{x^{0+1}}{0+1}}
    \onslide<+->{  + k}
    \\
    \onslide<+->{& = \frac{a}{15} (3x-1)^5}
    \onslide<+->{  - \frac{1}{c}\ln\abs{x}}
    \onslide<+->{  + \frac{x^{\sfrac{3}{2}}}{\sfrac{3}{2}}}
    \onslide<+->{  + 7x}
    \onslide<+->{  + k}
    \\
    \onslide<+->{& = \frac{a}{15} (3x-1)^5}
    \onslide<+->{  - \frac{1}{c}\ln\abs{x}}
    \onslide<+->{  + \frac{2x^{\sfrac{3}{2}}}{3}}
    \onslide<+->{  + 7x + k}
  \end{align*}
\end{resolution}

%------------------------------------------------------------------------------%

%------------------------------------------------------------------------------%
\section{Lista Mínima}

%------------------------------------------------------------------------------%
\begin{frame}{Lista Mínima}
  Estudar a Seção 2.2 da Apostila

  Exercícios: 1a-f, 2, 3

  \vfill
  Atenção: A prova é baseada no livro, não nas apresentações
\end{frame}

%------------------------------------------------------------------------------%
\end{document}
%------------------------------------------------------------------------------%
